\documentclass[10pt,a4paper]{amsart}
\usepackage[margin=19mm]{geometry}
\usepackage{amsmath,amssymb,mathtools}
\usepackage[scr=rsfs,cal=euler]{mathalfa}
\usepackage{xfrac}
\usepackage[unicode,hidelinks,bookmarks=false]{hyperref}
\newcommand{\ie}{i.e.\ }
\newcommand{\cf}{cf.\ }
\newcommand{\resp}{respectively\ }
\newcommand{\Subsection}{\subsection}
\providecommand{\nobreakdash}{\nobreak\hbox{-}}
\setlength{\emergencystretch}{2em}
\multlinegap0pt
\usepackage{bm}
\usepackage{bbm}
\usepackage{dsfont}
\usepackage{xspace}
\usepackage{cancel}
\usepackage{mathtools}
\usepackage{amsthm}
\makeatletter
\@ifundefined{thm}{\newtheorem{thm}{Theorem}}{}
\@ifundefined{lem}{\newtheorem{lem}[thm]{Lemma}}{}
\makeatother
\title[Fundamental solutions for parabolic equations and systems: u]{Fundamental solutions for parabolic equations and systems: universal existence, uniqueness, representation}
\author{Pascal Auscher, Khalid Baadi}
\date{Source: arXiv:2412.18436v3 (CC BY 4.0)}
\begin{document}
\maketitle

\noindent\emph{Source and licence.} Fundamental solutions for parabolic equations and systems: universal existence, uniqueness, representation. Version arXiv:2412.18436v3 (CC BY 4.0), distributed under the Creative Commons Attribution 4.0 International licence, as recorded in the arXiv licence metadata for this exact version and shown on the abstract page https://arxiv.org/abs/2412.18436v3. Full rights notice: licenses/arxiv-2412.18436-CC-BY-4.0.txt.\par\smallskip

\noindent\textit{Editorial scope.} Counted unit: the Lem whose source label is \texttt{lem:aprioriestimates} in the article (source0.tex lines 747--766, source label \texttt{lem:aprioriestimates}), delivered together with the article's own complete original proof of it (source0.tex lines 768--902, one \texttt{proof} environment of 135 lines). No mathematics is written, repaired or re-derived: statement and proof are the article's own text line for line. Required context retained uncounted: Spaces V\_alpha (lines 635--653); ISOMORPHISM (lines 681--683). Every definition, display and label of those lines is retained unchanged. Omitted from this package and not redistributed: 1--634, 654--680, 684--746, 767--767, 903--2091. Nothing inside a retained excerpt is omitted or altered: every retained body is delivered exactly as the article's lines are written, so no omission receipt of any kind accompanies this package. Citations: the retained excerpts carry no citation command at all, so no bibliography entry of the article is transcribed and no reference is redistributed. Reference adaptation: none was needed and none was made; every reference inside the delivered unit resolves inside it, so no unresolved reference or citation is produced. Printed slips kept literal: no printed word, symbol, hypothesis or number of the counted statement or of its proof was altered. Layout and class: the article's own class is \texttt{amsart} with packages including geometry, amsfonts, amsmath, amssymb, graphicx, caption, epstopdf, hyperref, xcolor, amsthm, float, pgfplots, listings, longtable; the wrapper is the lane's pinned \texttt{amsart} editorial wrapper at 10pt on A4, which restates the theorem-like environments the retained text needs and adds the article's own macro definitions that the retained text needs (none), with extra wrapper packages (bm, bbm, dsfont, xspace, cancel, mathtools). The delivered numbering is the unit's own. Assets: no retained span contains a figure environment, a graphics-inclusion command, a PDF-page inclusion, a table or a TikZ picture; no image file is copied into this package, the package declares no asset dependency and no image file is redistributed with it. Licence: version arXiv:2412.18436v3 of this article is distributed under the Creative Commons Attribution 4.0 International licence, as recorded in the arXiv licence metadata for this exact version and shown on the abstract page https://arxiv.org/abs/2412.18436v3. A full rights notice is delivered at licenses/arxiv-2412.18436-CC-BY-4.0.txt. Characters: every retained excerpt is pure ASCII, so the delivered engine, XeLaTeX with its default Latin Modern text font, reports no missing character.
\begin{lem}[Properties of intermediate spaces]\label{Spaces V_alpha} Fix  $-1 \leq \alpha \leq \alpha' \leq 1$. We have the following assertions. 
\begin{enumerate}
        \item $V_\alpha $ is a well-defined  subspace of  $\mathcal{S'}(\mathbb{R};E_\infty)$, $\left ( V_\alpha, \left \| \cdot \right \|_{V_\alpha} \right )$ is a Hilbert space, and we have 
                 \begin{equation*}
                     V_\alpha = \left \{ u \in \mathcal{S'}(\mathbb{R};E_\infty) : S^\alpha (S+\left | \tau \right |^{1/2})^{1-\alpha} \hat{u} \in L^2(\mathbb{R};H)\right \}, \ \left \| u \right \|_{V_\alpha} \sim  \| S^\alpha (S+\left | \tau \right |^{1/2})^{1-\alpha}\hat{u}  \|_{L^2(\mathbb{R};H)}.
                 \end{equation*}
        \item We have the following chain of continuous and dense inclusions:
        \begin{align*}
            \mathcal{S}(\mathbb{R};E_{-\infty}) \hookrightarrow V_\alpha \hookrightarrow V_{\alpha'} \hookrightarrow \mathcal{S'}(\mathbb{R};E_{\infty}).
        \end{align*}
         \item { $W_\alpha$ is a subspace of $\mathcal{S'}(\mathbb{R}; E_\infty)$, and $\left( W_\alpha, \left\| \cdot \right\|_{W_\alpha} \right)$ is a Hilbert space. We have a dense inclusion $\mathcal{S}_0(\mathbb{R}; E_{-\infty}) \hookrightarrow W_\alpha$, where
$\mathcal{S}_0(\mathbb{R}; E_{-\infty}) := \{ f \in \mathcal{S}(\mathbb{R}; E_{-\infty}) \mid \hat{f}(0) = 0 \}. $}
        \item Let $V_\alpha^\star$ denote the anti-dual space of $V_\alpha$ with respect to $\langle \cdot , \cdot \rangle_{L^2(\mathbb{R};H)}$. It is a subspace of $\mathcal{S}'(\mathbb{R};E_\infty)$ and  $V_\alpha^\star = L^2(\mathbb{R};D_{S,-1}) + W_{-\alpha}$ with the following estimate
    \begin{align*}     
   \left \| \omega  \right \|_{V_\alpha^\star} \sim \inf  \left\{  \left \| f \right \|_{L^2(\mathbb{R};H)}+\left \| g \right \|_{L^2(\mathbb{R};D_{S,-\alpha})}  \ : \ {\omega=Sf+D_t^{\frac{1-\alpha}{2}}g}\right\}.
\end{align*}
         
    \end{enumerate}
\end{lem}

\begin{thm}[Invertibility on intermediate spaces]\label{ISOMORPHISM}
    For all $\alpha\in [-1,1]$, the operator $ \partial_t + S^2 $, defined on $\mathcal{S}(\mathbb{R};E_{-\infty})$, extends to a bounded and invertible  operator $ A_\alpha :  V_\alpha \rightarrow V_{-\alpha}^\star$, which agrees with the restriction of $\partial_t + S^2$ acting on $\mathcal{S}'(\mathbb{R};E_\infty)$.
\end{thm}

\begin{lem}[A priori estimates for the Duhamel operator] \label{lem:aprioriestimates} 
Let $ f \in \mathcal{S}(\mathbb{R}; E_{-\infty}) $, and define $ u = Tf $. For the inequalities involving $ \|f\|_{W_\alpha} $, we additionally assume that $ f \in \mathcal{S}_0(\mathbb{R}; E_{-\infty})$. 
    
    \begin{enumerate} \item $u\in C_0(\mathbb{R}; H)$ and one has the following uniform bounds
        \begin{align*}\sup_{t \in \mathbb{R}} \left\| u(t)\right\|_H &\le \|f\|_{L^1(\mathbb{R};H)}\\
            \sup_{t \in \mathbb{R}} \left\| u(t)\right\|_H  &\le \frac{1}{\sqrt{2}} \|f \|_{L^2(\mathbb{R}; D_{S,-1})}\\
            \sup_{t \in \mathbb{R}} \left\| u(t)\right\|_H &\le C(\alpha) \|f\|_{W_\alpha}, \ \alpha\in [-1,0).
        \end{align*} 
        \item  $u \in L^2(\mathbb{R};D_{S,1})$ and one has the following energy inequalities
        \begin{align*}\left\|u \right\|_{L^2(\mathbb{R}; D_{S,1})} &\le \frac{1}{\sqrt{2}} \|f\|_{L^1(\mathbb{R};H)}\\
            \left\|u \right\|_{L^2(\mathbb{R}; D_{S,1})}  &\le  \|f \|_{L^2(\mathbb{R}; D_{S,-1})}\\
            \left\|u \right\|_{L^2(\mathbb{R}; D_{S,1})} &\le C'(\alpha) \|f\|_{W_\alpha}, \ \alpha\in [-1,1].
        \end{align*} 
        \item  $u \in V_\alpha$ for all $\alpha\in [-1,1]$ and one has the following bound
        \begin{align*}\|D_t^{\frac{1-\alpha}{2}}S^\alpha u\|_{L^2(\mathbb{R};H)}  \le  \|f\|_{W_\alpha}.
        \end{align*} 
         
            \end{enumerate}
        
\end{lem}

\begin{proof} That $u$ belongs to $L^\infty(\mathbb{R};H)$ with $\| u\|_{L^\infty(\mathbb{R};H)} \le \|f\|_{L^1(\mathbb{R};H)}$ has been already observed above. Note that $\partial_t u= f-T(S^2f)\in  L^\infty(\mathbb{R};H)$. Thus $u$ is Lipschitz, hence continuous. The limit 0 at $-\infty$ is clear from the fact that $\|f(s)\|_H$ has rapid decay and the contraction property of the semigroup.  As for the limit at $+\infty$, we write for fixed and large $A$ and $t>A$, 
$$ u(t)= e^{-(t-A)S^2}u(A)+ \int_A^t e^{-(t-s)S^2}f(s) \ \mathrm{d}s. $$
The first term tends to 0 in $H$ by properties of the semigroup and, for the second term, one uses again the contraction property and  rapid decay of $\|f(s)\|_H$. 

We are left with proving the remaining estimates. 

\

\paragraph{\textit{Step 1:  $\left\| u(t)\right\|_H \leq \frac{1}{\sqrt{2}} \left\|f \right\|_{L^2(\mathbb{R}; D_{S,-1})}$ for all $t \in \mathbb{R}$}}
Using Cauchy-Schwarz inequality, we have for all $t \in \mathbb{R}$ and $a \in H$, 
\begin{equation*}
    \int_{-\infty}^{t} | \langle  S^{-1}f(s) ,Se^{-(t-s)S^2} a \rangle_H | \ \mathrm d s \leq \frac{1}{\sqrt{2}}\left ( \int_{-\infty}^{t} \|S^{-1}f(s) \|^2_H \ \mathrm{d}s \right )^{1/2} \left\| a\right\|_H,
\end{equation*}
where we have used the quadratic equality 
\begin{equation*}
    \int_{0}^{\infty} \|Se^{-sS^2}a \|^2_H \mathrm d s =\frac{1}{2}\left\| a\right\|_H^2.
\end{equation*}
As 
\begin{equation*}
    \langle  u(t) , a \rangle_H = \int_{-\infty}^{t} \langle  S^{-1}f(s) , Se^{-(t-s)S^2}a \rangle_H \ \mathrm d s, 
\end{equation*}
we obtain the desired bound for $\|u(t)\|_H$.

\

\paragraph{\textit{Step 2:  $\left\|u \right\|_{L^2(\mathbb{R}; D_{S,1})} \le \frac{1}{\sqrt{2}} \|f\|_{L^1(\mathbb{R};H)}$}}
We observe that by Fubini's theorem, we have $\langle Su, \tilde f\rangle=  \langle u, S\tilde f\rangle=\langle f, \tilde u\rangle$, where $\tilde u=\tilde T(S\tilde f)$. Thus
\begin{equation*}
    |\langle Su, \tilde f\rangle| 
    \le  
    \left \| f  \right \|_{L^1(\mathbb{R};H)}  \| \tilde T(S\tilde f)   \|_{L^\infty(\mathbb{R};H)}
    \le 
    \frac{1}{\sqrt{2}} \left \| f  \right \|_{L^1(\mathbb{R};H)}  \| \tilde f   \|_{L^2(\mathbb{R};H)}
\end{equation*}
using step 1 for $\tilde T$. 

\

\paragraph{\textit{Step 3:  $\left\|u \right\|_{L^2(\mathbb{R}; D_{S,1})} \le  \|f\|_{L^2(\mathbb{R};D_{S,-1})}$}} We already know from step 2 that $\left\|u \right\|_{L^2(\mathbb{R}; D_{S,1})}$ is finite. To obtain the desired bound, we use again Fubini's theorem several times and obtain 
\begin{align*}
    \left\|Su \right\|_{L^2(\mathbb{R};H)}^2&
    =\int_{\mathbb{R}} \langle  Su(t) , Su(t) \rangle_H \ \mathrm{d} t 
    \\
    &
    = \int_{\mathbb{R}} \int_{-\infty}^{t}\int_{-\infty}^{t} \langle  S^2 e^{-(t-s)S^2}S^{-1}f(s) , S^2 e^{-(t-s')S^2}S^{-1}f(s')\rangle_H \ \mathrm{d} s  \mathrm{d} s'  \mathrm{d} t  \\ 
    &
    = \int_{\mathbb{R}} \int_{\mathbb{R}}\int_{\max(s,s')}^{+\infty} \langle  S^4 e^{-(2t-(s+s'))S^2}S^{-1}f(s) ,  S^{-1}f(s')\rangle_H \ \mathrm{d} t  \mathrm{d} s  \mathrm{d} s'\\
    & = \frac{1}{2} \int_{\mathbb{R}} \int_{\mathbb{R}} \langle  S^2 e^{-(2\max(s,s')-(s+s'))S^2}S^{-1}f(s) , S^{-1}f(s')\rangle_H  \ \mathrm{d} s  \mathrm{d} s'\\
    &=  \int_{\mathbb{R}} \int_{s\leq s'} \langle  S^2 e^{-(s'-s)S^2}S^{-1}f(s) , S^{-1}f(s')\rangle_H  \ \mathrm{d} s  \mathrm{d} s'\\
    &=  \int_{\mathbb{R}} \langle  Su(s') , S^{-1}f(s')\rangle_H \ \mathrm{d} s'.
\end{align*}
Using Cauchy-Schwarz inequality, we deduce that
\begin{align*}
     \left\|Su \right\|_{L^2(\mathbb{R};H)}^2 \leq  \left\|Su \right\|_{L^2(\mathbb{R};H)}  \left\|f \right\|_{L^2(\mathbb{R};D_{S,-1})}.
\end{align*}
Therefore 
\begin{align*}
     \left\|u \right\|_{L^2(\mathbb{R};D_{S,1})} \leq   \left\|f \right\|_{L^2(\mathbb{R};D_{S,-1})}.
\end{align*}

\

\paragraph{\textit{Step 4:  $\left\| u(t)\right\|_H \le C(\alpha) \|f\|_{W_\alpha}$, $ \alpha\in [-1,0)$}}
For all $a\in E_{-\infty}$, we define 
\begin{align*}
  \forall t \in \mathbb{R};\ \   \varphi_a(t):= \mathbb{1}_{(-\infty,0]}(t) e^{tS^2}a.
\end{align*}
Remark that when $a\in H$, $\varphi_a\in L^2(\mathbb{R}, D_{S,1})$, hence $\varphi_a \in L^2(\mathbb{R}, E_{-\infty})$.
 For  $t \in \mathbb{R}$,
\begin{align*}
    \langle u(t), a \rangle_H 
    &=  \int_{-\infty}^{0} \langle e ^{sS^2}f(s-t) , a \rangle_H \  \mathrm d s
    \\ 
    &= \int_{\mathbb{R}} \langle f(s-t), \varphi_a (s) \rangle_H \ \mathrm d s 
    \\ 
    &
    = \langle \tau_t g, D_t^{\frac{1+\alpha}{2}}S^{-\alpha}\varphi_a \rangle_{L^2(\mathbb{R};H),L^2(\mathbb{R};H)}. 
\end{align*}
In the calculation, we used $\hat f(0)=0$ (see Lemma \ref{Spaces V_alpha}, point (3)) and wrote $f=D_t^{\frac{1+\alpha}{2}}S^{-\alpha}g$ with $g\in L^2(\mathbb{R}; H)$, defined $\tau_tg(s)=g(s-t)$ and  used that translations commute with $D_t$.  If we show  that 
\begin{align*}\label{Hinfini estimation}
     \| D_t^{\frac{1+\alpha}{2}} S^{-\alpha}\varphi_a  \|_{L^2(\mathbb{R}; H)} = C(\alpha) \left \| a \right \|_H,
\end{align*} then
\begin{align*}
    \left | \langle u(t), a \rangle_H \right |\leq C(\alpha) \left \| g \right \|_{L^2(\mathbb{R}; H)}\left \| a \right \|_H,
\end{align*}
and  we may conclude using the density of $E_{-\infty}$ in $H$. To see this, 
applying Fourier transform to $\varphi_a$, we get 
\begin{align*}
    \forall \tau \in \mathbb{R}; \ \ \hat{\varphi_a}(\tau)= (-i \tau+S^2)^{-1}a, 
\end{align*}
so that
\begin{equation*}
    \left | \tau  \right |^{1+\alpha} \left \| S^{-\alpha}(-{i} \tau+S^2 )^{-1}a \right \|^2_H= |\tau|^{-1} \langle \psi(|\tau|^{-1/2}S) a, a \rangle_H
\end{equation*}
with $\psi(t)= t^{-2\alpha}(1+t^4)^{-1}$.  Using simple computations and Calder\'on's identity $$\int_{-\infty}
^\infty \psi(|\tau|^{-1/2}S)a \ \frac{\mathrm{d}\tau}{|\tau|}=  \int_0^\infty \frac{ t ^{-2\alpha }}{1+t^4} \frac{\mathrm d t}{ t }\  a,$$ we obtain 
\begin{align*}
    \int_{\mathbb{R}} \left | \tau  \right |^{1+\alpha} \left \|S^{-\alpha} \hat{\varphi_a }(\tau) \right \|^2_{H} \ \mathrm{d} \tau =   \int_0^\infty \frac{ t ^{-2\alpha }}{1+t^4} \frac{\mathrm d t}{ t }\ \left \| a \right \|^2_H,
\end{align*}
and conclude using Plancherel identity that $C(\alpha)^2= \frac{1}{2\pi}\int_0^\infty \frac{ t ^{-2\alpha }}{1+t^4} \frac{\mathrm d t}{ t }$.
\

\paragraph{\textit{Step 5:  $\left\|u \right\|_{L^2(\mathbb{R}; D_{S,1})}\le C'(\alpha) \|f\|_{W_{\alpha}}$, $ \alpha\in [-1,1]$}}
Since $f\in \mathcal{S}(\mathbb{R}; E_{-\infty})$, we know a priori that $u\in L^2(\mathbb{R}; D_{S,1})$ from Step 2 and $u$ agrees with the solution given by Fourier transform of Theorem \ref{ISOMORPHISM}. Hence, we can use Fourier transform to compute. We have 
\begin{align*}
    \hat{u}(\tau) = ({i}\tau +S^2)^{-1} \left | \tau  \right |^{\frac{1+\alpha}{2} }S^{-\alpha}g(\tau).
\end{align*}
where $f=D_t^{\frac{1+\alpha}{2}}S^{-\alpha}g$ with $g\in L^2(\mathbb{R}; H)$.
Hence
\begin{align*}
    \|S\hat u(\tau)\|^2_H
&= \langle (\tau^2+S^4)^{-1} |\tau|^{1+\alpha}S^{2-2\alpha} \hat g(\tau), \hat g(\tau)
\rangle_H \\
&= \|(|\tau|^{-1/2}S)^{1-\alpha}(1+(|\tau|^{-1/2}S)^4)^{-1/2}\hat g(\tau)\|_H^2 
\\
&
\le C'(\alpha)^2 \|\hat g(\tau)\|_H
^2
\end{align*}
with $C'(\alpha)= \sup_{t>0} t^{1-\alpha}(1+t^4)^{-1/2}<\infty$ when $-1\le \alpha \le 1$.

\

\paragraph{\textit{Step 6:  $\|D_t^{\frac{1-\alpha}{2}}S^\alpha u\|_{L^2(\mathbb{R};H)}  \le  \|f\|_{W_{\alpha}}$, $ \alpha\in [-1,1]$}}

We proceed as in step 5, and compute
\begin{align*}
    \||\tau|^{\frac{1-\alpha}{2}}S^\alpha \hat u(\tau)\|^2_H
= \langle (\tau^2+S^4)^{-1} |\tau|^{2} \hat g(\tau), \hat g(\tau)
\rangle_H 
\le  \|\hat g(\tau)\|_H
^2.
\end{align*}
 The conclusion follows. 
 \end{proof}

\end{document}
